%! Function Model Construction and Application — Unit 1, Lesson 1.14 — Exit Ticket (Answer Key)
\documentclass[10pt]{article}
\usepackage{precalculus-article}
\usepackage{precalculus-key}   % pulls in -boxes; do NOT also load -boxes
\newcommand{\TallMath}[1]{$\displaystyle #1\rule[-1.4em]{0pt}{3.2em}$}

\begin{document}

\pageheader{Unit 1, Lesson 1.14}{Exit Ticket --- Answer Key}
\namedateperiod

\vspace{0.12in}

\begin{enumerate}[itemsep=18pt]

\item A quadratic model gives the height of a projectile in feet:
      $f(x) = -2x^2 + 12x$, where $x$ is time in seconds.

      \begin{enumerate}[label=(\alph*), itemsep=10pt]
        \item Evaluate $f(3)$ and $f(5)$.

              \vspace{0.06in}
              $f(3) = -2({\color{keyred}\mathbf{3}})^2 + 12({\color{keyred}\mathbf{3}}) = {}$ \ans{18} ft

              \vspace{0.06in}
              $f(5) = -2({\color{keyred}\mathbf{5}})^2 + 12({\color{keyred}\mathbf{5}}) = {}$ \ans{10} ft

        \item Compute the average rate of change of $f$ on $[3, 5]$. Include units.

              \vspace{0.06in}
              $\text{ARC} = \dfrac{f(5) - f(3)}{5 - 3} = \dfrac{{\color{keyred}\mathbf{-8}}}{{\color{keyred}\mathbf{2}}} = {}$ \ans{$-4$}

              \vspace{0.04in}
              Units: \ans{feet per second (ft/s)}
      \end{enumerate}

\item A rational model for the force between two magnets is $P(d) = \dfrac{50}{d^2}$,
      where $d$ is the distance in meters between the magnets.

      State the domain restriction and explain in one sentence why it applies.

      \vspace{0.06in}
      Domain: \ans{$d > 0$} \quad Explanation: \ansline{Distance between physical objects must be a positive number; $d = 0$ would mean the magnets occupy the same point, which is physically impossible and causes division by zero.}

\item The average rate of change of the pen area $A(w)$ on $[6, 8]$ is $-4$ ft$^2$/ft.

      What does this tell you about the pen design on that interval? Answer in a complete sentence.

      \vspace{0.08in}
      \ansline{On average, each additional foot of width between 6 and 8 feet decreases the area by 4 square feet, meaning the pen is becoming less efficient (the design is past its optimal width).}

\end{enumerate}

\begin{teachernote}
Exit ticket item (1b): the negative ARC ($-4$ ft/s) means the projectile is losing height on
average between $t=3$ and $t=5$ seconds, consistent with the parabola opening downward.
Item (3) checks conceptual understanding of the sign and units of ARC.
Collect and sort: students missing units or sign need re-teaching of ARC interpretation.
\end{teachernote}

\end{document}
